\section{Introduction}\label{sec:introduction}

The tautological ring of the moduli space of stable curves is generated
by cotangent and kappa classes, together with the classes obtained by
forgetting markings and gluing nodes. Describing all relations among
these natural classes is a central problem in the intersection theory
of moduli spaces. Pixton proposed a uniform system of stable-graph
relations and conjectured that it is complete. We construct a relation
outside that system: an alternating product of corrected small diagonals
that vanishes in rational Chow, and therefore in rational cohomology.

\subsection{The relation problem}

Mumford initiated the systematic study of natural classes in the Chow
ring of the moduli space of curves \cite[Sections~4--5]{Mumford1983}.
Faber's calculations and geometric constructions on powers of the
universal curve led to a proposed description of the tautological ring
of the smooth unpointed space \(\mathcal M_g\), including its Gorenstein
property \cite[Conjectures~1--2]{Faber1999}. Faber and Zagier later gave
explicit hypergeometric formulas for relations among the kappa classes.
Pandharipande--Pixton proved these relations in Chow using stable
quotients and formulated their completeness on \(\mathcal M_g\)
\cite[Theorem~1 and Conjecture~2]{PandharipandePixton2021}.

On the compactification, a presentation must also account for boundary
strata and their intersections. Arbarello--Cornalba and
Graber--Pandharipande developed the forgetful and boundary calculus
\cite[Section~1]{ArbarelloCornalba1996},
\cite[Appendix~A]{GraberPandharipande2003}; in genus zero, Keel gave a
presentation of the entire Chow ring by boundary divisors
\cite[Section~4, Theorem~1]{Keel1992}.
Pixton's 2012 graph formula extends the Faber--Zagier system to
\(\Mbar_{g,n}\) and proposes a complete set of relations
\cite[Sections~1--3, Conjecture~2]{Pixton}.
Pandharipande--Pixton--Zvonkine proved the cohomological validity of this
system using the \(3\)-spin CohFT \cite{PPZ}; Janda proved its Chow
validity using equivariant Gromov--Witten theory of \(\mathbb P^1\)
\cite{Janda}. These results establish relations, while completeness asks
whether any further relations exist.

We fix the original relation space precisely. Let \(\Salg_{g,n}\) be the
strata algebra of formal decorated stable graphs, with decoration degree
truncated at the dimension of each vertex factor. Its realization is
\[
q_{g,n}\colon\Salg_{g,n}\longrightarrow\CH^*(\Mbar_{g,n};\Q),
\qquad [\Gamma,\gamma]\longmapsto\xi_{\Gamma*}\gamma.
\]
The gluing pushforward generators are unnormalized, as in
\cite[Section~0.3]{PPZ}. Let \(\Pspan_{g,n}\) be the span obtained by
placing one of Pixton's graph-formula relations at a vertex of an outer
stable graph, placing arbitrary strata classes at the other vertices,
and gluing. Pixton's Proposition~1 shows that this span is already an
ideal, stable under the gluing and forgetting pushforwards and pullbacks
\cite{Pixton}. Thus it includes these formal consequences from the
outset. The smaller \(0/1\)-marking seed system recovers the same span
under the operations described in \cite[Section~3.5]{PPZ}.
Section~\ref{sec:strata} gives the full recipe and all conventions.

Writing \(R^*(\Mbar_{g,n})=\im q_{g,n}\), the Chow completeness question is
whether
\[
\ker q_{g,n}=\Pspan_{g,n}
\]
for every stable pair. The known validity theorem gives the inclusion
\(\Pspan_{g,n}\subseteq\ker q_{g,n}\). Our example makes it strict.

\subsection{The counterexample}

The construction uses \(g\) colors and one three-label factor for each
increasing triple of colors. Each factor is a product of two
rational-tail divisor sums, corrected by a class supported on two-node
rational bridges. We multiply these factors and alternate over
relabelings that preserve the color of each slot.
Definition~\ref{def:candidate} specifies the resulting formal class \(Y\).

\begin{theorem}\label{thm:main}
Set
\[
g=10^{60},\qquad D_*=\binom g3,\qquad n=3D_*.
\]
The explicitly defined class \(Y\in \Salg_{g,n}^{2D_*}\) of
Definition~\ref{def:candidate} satisfies
\[
q_{g,n}(Y)=0,\qquad Y\notin\Pspan_{g,n}.
\]
Consequently Pixton's original relation span does not exhaust the
tautological Chow relations: the equality
\(\ker q_{g,n}=\Pspan_{g,n}\) fails for this stable pair.
\end{theorem}

The same class gives the rational-cohomological consequence. To specify
that formulation, write \(q^{\mathrm{CH}}_{g,n}=q_{g,n}\) and define, in
formal degree \(d\),
\[
q^H_{g,n}=\operatorname{cl}\circ q^{\mathrm{CH}}_{g,n}
\colon\Salg^d_{g,n}\longrightarrow
H^{2d}\!\left(\Mbar_{g,n}(\C);\Q\right),
\]
where \(\operatorname{cl}\) is the rational cycle-class map for the complex
Deligne--Mumford stack. Its image is the cohomological tautological ring
\(RH^*=\operatorname{cl}(R^*)\). Completeness in both \(R^*\) and \(RH^*\)
is explicitly recorded as Pixton's conjecture in
\cite[Section~1, Conjecture~1]{Canning2025}.

\begin{corollary}[Rational-cohomological incompleteness]
\label{cor:cohomological-incompleteness}
For the stable pair and class \(Y\) of Theorem~\ref{thm:main},
\[
q^H_{g,n}(Y)=0
\quad\text{in }H^{4D_*}\!\left(\Mbar_{g,n}(\C);\Q\right),
\qquad
Y\notin\Pspan_{g,n}.
\]
In particular, \(\Pspan_{g,n}\subsetneq\ker q^H_{g,n}\).
\end{corollary}

\begin{proof}
Cycle class commutes with Chern classes, products, and proper
pushforward. In particular,
\(\operatorname{cl}(\xi_{\Gamma*}\gamma)
=\xi_{\Gamma*}\operatorname{cl}(\gamma)\), and the same compatibility
applies to the forgetting-pushforward definition of kappa classes.
Thus \(q^H_{g,n}\) is the cohomological graph realization of
\cite[Section~0.3]{PPZ}, with the same vertexwise truncation and
unnormalized generators. Applying cycle class to
Theorem~\ref{thm:main} gives \(q^H_{g,n}(Y)=0\); nonmembership remains
an assertion in the same formal strata algebra. The strict inclusion
follows from
\[
\Pspan_{g,n}\subseteq\ker q^{\mathrm{CH}}_{g,n}
\subseteq\ker q^H_{g,n}.
\]
Passing from Chow vanishing to cohomological vanishing uses the second
inclusion; it requires no injectivity of cycle class.
\end{proof}

The theorem and corollary give negative answers to the Chow and
rational-cohomological completeness questions for the full original
system. The large parameters leave ample room in the support and color
estimates; they are not proposed as minimal parameters.

Completeness is distinct from the Gorenstein property. The latter asks
for nondegenerate multiplication pairings in complementary degrees,
whereas completeness asks for all relations among the generators.
Pixton proved that \(R^*(\mathcal M_{25})\) is not Gorenstein while
establishing Faber--Zagier completeness there
\cite[Theorem~1.1 and Section~2.2]{Pixton2026}.
Canning--Larson--Schmitt proved completeness together with
non-Gorenstein behavior on the compact-type spaces
\(\mathcal M^{\mathrm{ct}}_6\), \(\mathcal M^{\mathrm{ct}}_{5,2}\), and
\(\mathcal M^{\mathrm{ct}}_7\)
\cite[Theorems~1 and~7]{CanningLarsonSchmitt2026}.
Our relation lies on the full stable compactification.

There is also a distinction between completeness and comparison of
methods for producing relations. Janda's comparison theorem identifies
the systems arising from cancellation of reconstruction singularities
for convergent semisimple CohFTs with nonempty discriminant
\cite[Definition~3.3.1 and Theorem~3.3.6]{JandaDiscriminant}.
It concerns the relations generated by those reconstructions; it does
not identify the entire Chow relation kernel. The comparison throughout
this paper is with the explicitly fixed span \(\Pspan_{g,n}\).

\subsection{Proof and methods}

The proof compares two algebras of cubic coefficients. A formal
super Frobenius algebra can carry independent odd cubic coefficients
while satisfying the identities needed to annihilate the original
Pixton span. The actual Chow correspondences of a curve satisfy an
additional nilpotence identity obtained from a surface Quot scheme.
The first algebra proves nonmembership; the second proves vanishing.
Figure~\ref{fig:proof-map} separates these two assertions about \(Y\).

\paragraph{The formal test.}
Section~\ref{sec:detector} constructs a linear map on decorated graphs
by discarding graphs outside the rational-tails locus and evaluating
rational branches by Frobenius multiplication and coproduct. This uses
the graph calculus familiar from two-dimensional topological field
theory \cite[Section~2 and Theorem~3]{Abrams1996}.
The odd symplectic curve pairing and the vanishing of contractions
between isotropically supported corrections have precedents in
Pandharipande--Zvonkine's CohFT constructions
\cite[Sections~1.1--1.5 and Remark~11]{PZ2019}.
Here the coefficient ring is the exterior algebra on independent
\(\theta_{ijk}\), one for every color triple; we verify the graph
identities directly. Every element of \(\Pspan_{g,n}\) has zero test,
while the test of \(Y\) has a nonzero coefficient proportional to
\(\prod_{i<j<k}\theta_{ijk}\).

\paragraph{The smooth-family Chow identity.}
Section~\ref{sec:tensors} splits the correspondence object of a smooth
curve into two even summands and an odd symplectic summand. The
coefficient-algebra construction is a concrete form of O'Sullivan's
semisimple tensor reconstruction
\cite[Lemma~2.3.1]{OSullivan2011}. It retains all actual relative Chow
correspondences, so its faithfulness is an assertion about Chow, not
about cohomological realization. The small diagonal projected to three
odd summands gives an alternating cubic \(\varphi\) with odd
coefficients. Over a pointed curve this is the modified diagonal of
Gross--Schoen \cite[Equation~(0.1) and Sections~2--3]{GrossSchoen1995};
the family construction also retains the necessary base corrections.

Sections~\ref{sec:quot}--\ref{sec:coefficient} prove the additional
nilpotence. The surface is the square of the curve. A relative
Quot scheme with fixed determinant produces a relation for products
of the cubic on isotropic vector sets. Its obstruction theory uses
derived perfect-complex moduli and determinant maps \cite{TV,STV},
virtual pullbacks \cite{BF,Manolache}, and virtual localization
\cite{GP}. Surface Quot localization and relative Chow identities
appear in Marian--Oprea--Pandharipande \cite{MOPSegre}; the point
obstruction factors have the antecedent in
\cite[Lemma~34 and Section~5.2]{OPQuot}.
Here the calculation keeps the determinant constraint, arbitrary
irregularity, and every marked surface output.

The nested-Hilbert recursion of Ellingsrud--G\"ottsche--Lehn
\cite{EGL} reduces that marked calculation to decorated equality
diagonals. Kong develops relative forms of this recursion and, for
curve families, formulas retaining one and two marked outputs
\cite[Section~2.1 and Lemma~3.4]{Kong2026}.
We prove the version retaining an arbitrary finite set of surface
outputs. Only bare triple blocks survive the relevant codimension
bound. Section~\ref{sec:coefficient} computes their common scalar and
proves it is nonzero, using symmetric-power and residue methods related
to \cite[Equation~(9), Theorem~2, and Lemma~3]{MOPCombinatorics}.
It then converts the isotropic identity into vanishing tests for
sufficiently long products with arbitrary color inputs
(Lemma~\ref{lem:large-vertex-test}). Finite faithful polarization returns
a Chow identity modulo classes supported over the prescribed positive
codimension in the base.

\paragraph{Passage to the compactification.}
Section~\ref{sec:compactification} compares the corrected stable cubic
with these projected smooth-curve cubics. The framework uses Hassett's
weighted stable curves \cite{Hassett} and configuration-space blowups
\cite{FultonMacPherson,Routis}; Petersen's proof of Tavakol's
presentation conjecture treats the corresponding rational-tails
comparison over a smooth core \cite{TavakolRT,PetersenWonderful}.
The nodal comparison here requires a rational-bridge correction with
coefficient \(-1/2\). A three-position calculation shows that it
cancels the remaining triple-exceptional contribution. The errors are
then either a two-slot operation times a one-slot operation from a
bounded span, or operations supported over positive base codimension.

Section~\ref{sec:support} uses these errors to lower the dimension of
the image of the supporting cycle in the base. Chow-localization
support methods occur in Graber--Vakil's filtration by rational
components \cite[Theorem~1.1 and Section~4.2]{GraberVakil}; the induction
here acts on arbitrary cycle multipliers over the base parametrizing
the curve components.
When one component carries almost all the genus, the surface identity
advances support. Otherwise a repeated-color identity supplies a
smaller vanishing factor, using the curve algebra centered at a marked
node branch,
Schur--Weyl alternation, and symplectic contraction. The final color
budget leaves enough entries for every advancement, until the support
is empty. Pullback from the weighted space then proves \(q_{g,n}(Y)=0\).

Thus Sections~\ref{sec:strata}--\ref{sec:detector} construct and detect
the formal candidate; Sections~\ref{sec:tensors}--\ref{sec:coefficient}
prove the smooth-family identity; and
Sections~\ref{sec:compactification}--\ref{sec:support} extend its
vanishing to the compactification. Table~\ref{tab:notation} collects
the notation used across these parts.

\begin{figure}[tb]
\centering
\begin{tikzpicture}[
 box/.style={draw,rounded corners=2pt,align=center,text width=5.7cm,
             minimum height=1.0cm,inner sep=5pt,font=\small},
 result/.style={box,fill=black!5},
 flow/.style={-{Stealth[length=2mm]},semithick},
 node distance=6mm]
\node[box] (candidate) {Corrected three-slot factors\\ and their alternation \(Y\)\\
 Definition~\ref{def:candidate}};
\node[box,below=8mm of candidate,xshift=-3.2cm] (formal) {Formal Frobenius algebra\\ over an exterior ring\\
 Rational-tail detector (Section~\ref{sec:detector})};
\node[box,below=8mm of candidate,xshift=3.2cm] (chow) {Faithful model of relative Chow correspondences\\
 Finite color tests (Section~\ref{sec:tensors})};
\node[result,below=8mm of formal] (outside) {\(t_n(\Pspan_{g,n})=0\), \(t_n(Y)\ne0\)\\
 Hence \(Y\notin\Pspan_{g,n}\)};
\node[box,below=8mm of chow] (surface) {Marked Quot relation \(aT_\gamma^{t_0}=0\)\\
 Section~\ref{sec:quot}; \(a\ne0\) in Section~\ref{sec:coefficient}};
\node[box,below=8mm of surface] (colors) {Arbitrary color products\\
 Lemma~\ref{lem:large-vertex-test} (Section~\ref{sec:coefficient})};
\node[box,below=8mm of colors] (boundary) {Controlled compactification\\ and support induction\\
 Sections~\ref{sec:compactification}--\ref{sec:support}};
\node[result,below=8mm of boundary] (zero) {\(q_{g,n}(Y)=0\) in rational Chow};
\draw[flow] (candidate.south) -- ++(0,-.25) -| (formal.north);
\draw[flow] (candidate.south) -- ++(0,-.25) -| (chow.north);
\draw[flow] (formal) -- (outside);
\draw[flow] (chow) -- (surface);
\draw[flow] (surface) -- (colors);
\draw[flow] (colors) -- (boundary);
\draw[flow] (boundary) -- (zero);
\end{tikzpicture}
\caption{The two assertions about the same class \(Y\). The formal
branch is independent of Chow vanishing. In the geometric branch,
the nonzero scalar completes the isotropic surface bound. The
exterior-ideal argument then gives arbitrary-color vanishing modulo base
support before the passage to the compactification.}
\label{fig:proof-map}
\end{figure}


\subsection{Conventions}

All varieties, stacks, and families are over \(\C\), and Chow groups
have rational coefficients. On a smooth space we identify Chow classes
with operational classes when pulling them to a test space.
Products of stacks are understood in rational Chow theory
\cite[Theorem~2.1.12]{Kresch}.
We use proper rational pushforward for relative Deligne--Mumford maps
and its compatibilities from \cite[Theorem~B.17 and
Propositions~B.18--B.19]{BaeSchmitt}. Smooth Deligne--Mumford stacks have
the operational identification in \cite[Theorem~C.6 and
Corollary~C.9]{BaeSchmitt}; for smooth quotient stacks, see also
\cite[Propositions~4 and~19]{EdidinGraham}.
Every formal series calculation is
truncated at a specified coefficient, so only finitely many terms occur.

Parity is distinct from Chow codimension. In a super tensor product,
interchanging homogeneous factors of parities \(p,q\) gives
\((-1)^{pq}\). In particular, on an odd vector space the
super-symmetric tensors are ordinary alternating tensors.
All finite averaging and polarization arguments take place in
characteristic zero.


\begin{table}[tb]
\centering
\small
\begin{tabular}{@{}p{.20\textwidth}p{.53\textwidth}p{.19\textwidth}@{}}
\toprule
Notation & Meaning & Definition\\
\midrule
\(\Salg_{g,n},q_{g,n},\Pspan_{g,n}\) & Formal strata algebra, Chow realization, original Pixton span & Section~\ref{sec:strata}\\
\(D_*,n,Y\) & Number of cubic factors, number of labels, candidate relation & Definition~\ref{def:candidate}\\
\(R,V,t_m\) & Exterior coefficient ring, formal Frobenius algebra, detector & Section~\ref{sec:detector}\\
\(e'_C,\eta,H_C\) & Corrected degree-one class, pure projector, complementary curve object & Section~\ref{sec:tensors}\\
\(\mathfrak B,W_b,\varphi\) & Faithful Chow coefficients, odd symplectic space, projected cubic & Proposition~\ref{prop:tensor-test}\\
\(t_0,T_\gamma,a_{N,t_0}\) & Nilpotence exponent, isotropic cubic contraction, localization scalar & Sections~\ref{sec:quot}--\ref{sec:coefficient}\\
\(F^0_{h,B},G_{h,B}\) & Corrected weighted cubic and its chosen pure lift & Section~\ref{sec:compactification}\\
\(X^\Gamma,\mathcal B_\Gamma,\mathcal S_\Gamma\) & Cut chart, compact heavy base, smooth-core open & Section~\ref{sec:compactification}\\
\bottomrule
\end{tabular}
\caption{Notation used across the proof. The formal coefficient ring
\(R\) and the faithful Chow coefficient algebra \(\mathfrak B\) have
different roles.}
\label{tab:notation}
\end{table}
